\unnumberedchapter{Project Planning Guide}\label{ch:project-planning-guide}

\chaptersubtitle{A structured documentation template for new projects.
Technology-agnostic: independent of any specific framework or language.}

\chapterrule

\unnumberedsection{true}{What This Guide Is}\label{project-planning-guide:what-this-guide-is}

This is a structured documentation template to complete before~--- or at the very start of~--- any new project. Every section a project needs~--- across problem framing, architecture, code craft, operations, and delivery~--- is defined here at the conceptual level, ready to be filled in with project-specific detail.

The Guide is \textbf{technology-agnostic}. It does not prescribe which language to write in, which database to use, or which platform to deploy on. It defines \emph{what topics must be documented}, \emph{what each section should contain}, and \emph{what a well-completed entry looks like}.

The structure mirrors what real project documentation actually contains~--- the sections any project's docs would cover are the sections this Guide prompts you to complete.

\unnumberedsection{true}{How to Read This Guide}\label{project-planning-guide:how-to-read-this-guide}

For each subsection:

\begin{itemize}
\tightlist
\item
  \textbf{One-line definition}~--- what this section \emph{is}
\item
  \textbf{Why it matters}~--- what the project loses if this section is skipped or left vague
\item
  \textbf{What this topic covers}~--- the specific questions and details to address
\item
  \textbf{You understand this when}~--- a completeness check: the standard a well-filled section should meet
\end{itemize}

\noindent{}This is not a tutorial. It is a documentation prompt. The entries you write for each section are the deliverable; this Guide is the structure that tells you what to write.

\unnumberedsection{true}{Section 1 --- Problem Statement \& Goals}\label{project-planning-guide:section-1--problem-statement--goals}

\subsection*{1.1 Background \& Context}\phantomsection\label{project-planning-guide:11-background--context}

The history and surrounding conditions that gave rise to the project~--- the situation, team, system, or workflow this work plugs into.

\textbf{Why it matters:} Without context, a project's purpose is ambiguous. Engineers without context build technically correct systems that solve the wrong problem.

\textbf{What this topic covers:}

\begin{itemize}
\tightlist
\item
  The triggering event or pain point that initiated this project
\item
  Upstream and downstream systems
\item
  The team, process, or workflow this fits into
\item
  Previous attempts and why they didn't work
\item
  The history that informs current constraints
\end{itemize}

\noindent{}\textbf{You understand this when} you can explain in two paragraphs why this project exists \emph{now}, what alternatives were considered, and what would happen if it weren't done.

\subsection*{1.2 Core Problem}\phantomsection\label{project-planning-guide:12-core-problem}

The specific, concrete problem this project solves~--- narrow enough to be solvable, broad enough to matter.

\textbf{Why it matters:} A vague problem statement allows scope creep and produces a system that satisfies no one. A precise problem statement is the foundation for every later decision.

\textbf{What this topic covers:}

\begin{itemize}
\tightlist
\item
  The current pain, stated in 1--2 sentences
\item
  How often the problem occurs
\item
  Who is affected and what it costs them
\item
  The root cause where known
\item
  How to distinguish the problem from any proposed solution
\end{itemize}

\noindent{}\textbf{You understand this when} you can describe the problem in 30 seconds and a stranger understands what's broken without needing technical context.

\subsection*{1.3 Goals \& Success Criteria}\phantomsection\label{project-planning-guide:13-goals--success-criteria}

What ``done'' looks like~--- observable, ideally measurable conditions that indicate the problem is solved.

\textbf{Why it matters:} Goals like ``improve performance'' never tell you when to stop. Goals like ``reduce p99 latency from 800 ms to 200 ms'' do.

\textbf{What this topic covers:}

\begin{itemize}
\tightlist
\item
  Each goal as a specific outcome
\item
  Each goal paired with an observable success condition
\item
  The difference between vague aspirations and measurable criteria
\item
  Functional vs non-functional goals
\item
  Goals measurable before, during, and after the work
\end{itemize}

\noindent{}\textbf{You understand this when} every goal has a clear, testable success condition that even an outsider could verify.

\subsection*{1.4 Out of Scope}\phantomsection\label{project-planning-guide:14-out-of-scope}

What this project explicitly does \textbf{not} address~--- the boundaries that protect against scope creep.

\textbf{Why it matters:} Without explicit scope boundaries, every stakeholder assumes their pet feature is in scope. Listing what's out makes it possible to ship.

\textbf{What this topic covers:}

\begin{itemize}
\tightlist
\item
  Specific features or capabilities not being built
\item
  The reason each thing is excluded
\item
  Items deferred to a future phase, with the reason
\item
  The distinction between ``not yet'' and ``never''
\item
  How to negotiate scope when new requests arrive
\end{itemize}

\noindent{}\textbf{You understand this when} the team can answer ``is X in scope?'' without re-litigating the question every meeting.

\subsection*{1.5 Stakeholders \& Users}\phantomsection\label{project-planning-guide:15-stakeholders--users}

Who is affected by this project~--- users, beneficiaries, owners, and decision-makers.

\textbf{Why it matters:} Different roles need different things. Confusing ``user'' with ``decision-maker'' or ``owner'' with ``beneficiary'' causes communication failures and missed requirements.

\textbf{What this topic covers:}

\begin{itemize}
\tightlist
\item
  Primary users who interact day-to-day
\item
  Beneficiaries who consume output indirectly
\item
  Owners and maintainers responsible for the running system
\item
  Stakeholders with input on priorities or compliance
\item
  How to talk to each group about the work
\end{itemize}

\noindent{}\textbf{You understand this when} you can name every role, what they care about, how often they're consulted, and who has final say.

\subsection*{1.6 Constraints \& Assumptions}\phantomsection\label{project-planning-guide:16-constraints--assumptions}

Hard constraints (non-negotiable) and assumptions (things assumed true, that would change the design if false).

\textbf{Why it matters:} Implicit constraints become bugs. Implicit assumptions become outages. Both must be made explicit before they bite.

\textbf{What this topic covers:}

\begin{itemize}
\tightlist
\item
  Technical constraints (platforms, languages, integrations)
\item
  Legal and compliance constraints
\item
  Budget and timeline constraints
\item
  Team and capability constraints
\item
  Assumptions about user behavior, data, and environment
\item
  Distinguishing constraints (hard) from assumptions (probabilistic)
\end{itemize}

\noindent{}\textbf{You understand this when} an outside engineer reading the constraints list can predict where the design is going to bend before they see it.

\unnumberedsection{true}{Section 2 --- Architecture \& Design}\label{project-planning-guide:section-2--architecture--design}

\subsection*{2.1 System Overview}\phantomsection\label{project-planning-guide:21-system-overview}

A narrative of three to five steps describing what happens from trigger to output, in plain language, including how the system runs as a process.

\textbf{Why it matters:} Without an overview, every other architecture decision is unmoored. The overview is the elevator pitch for the whole system.

\textbf{What this topic covers:}

\begin{itemize}
\tightlist
\item
  Entry points~--- what initiates the system (user action, cron, event, API call)
\item
  Process model~--- how the application exists as a running process (single process, worker pool, scheduled job) and what supervises it
\item
  Core processing~--- what the system does with its input
\item
  Output~--- what the system produces and where it goes
\item
  Runtime lifecycle~--- startup → ready → handle → shutdown
\end{itemize}

\noindent{}\textbf{You understand this when} you can sketch the system for someone in five sentences, with no implementation jargon~--- the details belong in the sections that follow.

\subsection*{2.2 Architecture Diagram}\phantomsection\label{project-planning-guide:22-architecture-diagram}

A visual of the system showing components, data flow, external systems, storage, and APIs.

\textbf{Why it matters:} A diagram often communicates faster than prose. A team without diagrams reinvents the architecture in every conversation.

\textbf{What this topic covers:}

\begin{itemize}
\tightlist
\item
  Diagram types~--- flowchart, data-flow, system context, sequence, deployment
\item
  What to show~--- components, arrows, external systems, storage, APIs
\item
  The C4 model (context, container, component, code) and when each level fits
\item
  ASCII diagrams as a viable medium when there's ``no time for diagrams''
\item
  The discipline of keeping diagrams current
\end{itemize}

\noindent{}\textbf{You understand this when} you can produce a diagram for a system you just understood, at the right level of detail for the audience.

\subsection*{2.3 Component Breakdown}\phantomsection\label{project-planning-guide:23-component-breakdown}

Each major component, its responsibility, its inputs and outputs, its dependencies.

\textbf{Why it matters:} A system is its components. Without explicit responsibility boundaries, ownership is unclear and changes cascade unpredictably.

\textbf{What this topic covers:}

\begin{itemize}
\tightlist
\item
  One responsibility per component
\item
  Inputs each component accepts
\item
  Outputs each component produces
\item
  Dependencies between components
\item
  Interfaces~--- what's exposed, what's internal
\item
  How to recognize when one component is doing too much (or too little)
\end{itemize}

\noindent{}\textbf{You understand this when} you can tell whether a change should live in component A or component B, and defend the choice.

\subsection*{2.4 Data Flow}\phantomsection\label{project-planning-guide:24-data-flow}

The journey data takes through the system, from source to final destination, including the end-to-end request lifecycle.

\textbf{Why it matters:} Bugs in distributed systems are usually data-flow bugs~--- data that didn't arrive, arrived in the wrong order, or arrived in the wrong form. Tracing data flow is one of the most effective debugging techniques.

\textbf{What this topic covers:}

\begin{itemize}
\tightlist
\item
  Source~--- where data originates
\item
  Ingestion~--- how data enters the system
\item
  Processing~--- transformations, validations, enrichments
\item
  Storage~--- intermediate and final results
\item
  Output~--- final form and destination
\item
  Error paths~--- what happens to data that fails validation or processing
\item
  The request lifecycle in detail~--- client → network → kernel → application → database → response → back
\end{itemize}

\noindent{}\textbf{You understand this when} you can trace any input through the system and predict every transformation it undergoes.

\subsection*{2.5 External Integrations}\phantomsection\label{project-planning-guide:25-external-integrations}

Every external system, API, or service the project depends on~--- and the contract with each one.

\textbf{Why it matters:} External integrations are where systems meet. A bad integration contract is harder to fix than internal bugs because it requires both sides to change.

\textbf{What this topic covers:}

\begin{itemize}
\tightlist
\item
  Each integration's purpose
\item
  Authentication method (API key, OAuth, mTLS, service account)
\item
  Failure modes and how the system degrades when the integration is unavailable
\item
  Contract dimensions~--- shape (REST/GraphQL/gRPC), versioning, identity, authorization, failure semantics
\item
  Idempotency for safely retryable operations
\item
  Rate limits imposed by external services
\item
  Webhook handling for incoming events
\end{itemize}

\noindent{}\textbf{You understand this when} for every external dependency, you can answer ``what happens when this is unavailable?'' and ``how do we handle version skew safely?''

\subsection*{2.6 Data Storage \& Schema}\phantomsection\label{project-planning-guide:26-data-storage--schema}

Every database, file store, and cache used~--- including the schema design and storage choice rationale.

\textbf{Why it matters:} Schema choices outlive everything else. The application can be rewritten; the database structure usually cannot.

\textbf{What this topic covers:}

\begin{itemize}
\tightlist
\item
  Storage type (relational, document, key-value, graph, columnar, time-series, search, object/blob)
\item
  Schema design~--- entities, attributes, relationships
\item
  Normalization (1NF--BCNF) and when to denormalize
\item
  Key strategies (auto-increment, UUID, ULID; surrogate vs natural)
\item
  Indexing strategy
\item
  Transactions and consistency model
\item
  Retention policy, archival, deletion behavior
\item
  Schema evolution patterns (expand/contract, dual-write, backfills)
\end{itemize}

\noindent{}\textbf{You understand this when} for every entity in the system, you can defend the choice of storage type, the schema design, and the indexing strategy.

\subsection*{2.7 Key Design Decisions}\phantomsection\label{project-planning-guide:27-key-design-decisions}

The significant decisions made during architecture and design~--- and crucially, \emph{why}.

\textbf{Why it matters:} Decisions made today are forgotten in six months. Without a written record, every architectural debate gets re-litigated.

\textbf{What this topic covers:}

\begin{itemize}
\tightlist
\item
  The decision itself
\item
  Alternatives considered
\item
  Rationale~--- why this option won
\item
  Trade-offs accepted
\item
  Date and owner of the decision
\item
  ADR (Architecture Decision Record) format and lifecycle
\item
  When a decision deserves an ADR and when it doesn't
\end{itemize}

\noindent{}\textbf{You understand this when} you can read an ADR and reconstruct the discussion that produced it, including the rejected alternatives.

\subsection*{2.8 Security Considerations}\phantomsection\label{project-planning-guide:28-security-considerations}

How the system protects against misuse, attack, and accidental exposure.

\textbf{Why it matters:} Security cannot be bolted on at the end. It must be designed in. A breach costs more than a hundred features.

\textbf{What this topic covers:}

\begin{itemize}
\tightlist
\item
  Authentication~--- how identity is verified (sessions, tokens, OAuth, mTLS)
\item
  Authorization~--- who can do what (RBAC, ABAC, ownership-based)
\item
  Secrets management~--- where credentials live, how they rotate
\item
  Data sensitivity~--- classifying data (PII, financial, public) and handling each appropriately
\item
  Transport security~--- HTTPS everywhere, TLS configuration
\item
  Input validation and injection prevention (SQL, command, XSS, SSRF)
\item
  Defense in depth~--- multiple layers, no single point of failure
\item
  Logging policy~--- what to log, what \emph{not} to log
\item
  OWASP Top 10 awareness
\end{itemize}

\noindent{}\textbf{You understand this when} you can threat-model a new feature in your head and identify the top three risks before implementation begins.

\unnumberedsection{true}{Section 3 --- Developer Guide}\label{project-planning-guide:section-3--developer-guide}

\subsection*{3.1 Prerequisites}\phantomsection\label{project-planning-guide:31-prerequisites}

Everything that must exist on a machine before the project can run.

\textbf{Why it matters:} Undocumented prerequisites are a leading cause of ``works on my machine'' failures.

\textbf{What this topic covers:}

\begin{itemize}
\tightlist
\item
  Runtime~--- language and version
\item
  Package manager
\item
  External tools and system libraries
\item
  Credentials required (env vars, API keys, certificates)
\item
  System requirements (OS, RAM, disk, network)
\item
  How to verify each prerequisite is present
\end{itemize}

\noindent{}\textbf{You understand this when} a colleague can install every documented prerequisite on a fresh machine without asking anyone for help.

\subsection*{3.2 Repository Structure}\phantomsection\label{project-planning-guide:32-repository-structure}

A map of the codebase~--- where things live, the conventions followed.

\textbf{Why it matters:} Developers should be able to find code by following the structure, not by searching.

\textbf{What this topic covers:}

\begin{itemize}
\tightlist
\item
  Top-level directories and what each contains
\item
  Application code organization
\item
  Where configuration, tests, scripts, documentation, and fixtures live
\item
  A single, explicit composition point where the application is assembled (in Flask, the application factory)
\item
  Naming conventions for files and directories
\end{itemize}

\noindent{}\textbf{You understand this when} you can predict where to add a new feature file before opening the project.

\subsection*{3.3 Environment Setup}\phantomsection\label{project-planning-guide:33-environment-setup}

The step-by-step process for getting the project running locally.

\textbf{Why it matters:} The setup process is the first impression. Bad setup means lost developers.

\textbf{What this topic covers:}

\begin{itemize}
\tightlist
\item
  Cloning the repository
\item
  Creating and activating an isolated environment
\item
  Installing dependencies
\item
  Copying and configuring the environment file
\item
  Running migrations or seed scripts
\item
  Verifying setup with a smoke test
\item
  Common setup failures and their fixes
\end{itemize}

\noindent{}\textbf{You understand this when} someone new can go from \texttt{git\ clone} to a running local instance in under 15 minutes following only your documented steps.

\subsection*{3.4 Running Locally}\phantomsection\label{project-planning-guide:34-running-locally}

How to start the application, in development mode and in production-like mode.

\textbf{Why it matters:} Running locally is where most code is written and debugged. Subtle differences between local and production cause bugs that only surface late.

\textbf{What this topic covers:}

\begin{itemize}
\tightlist
\item
  Start command for development mode
\item
  Start command for production-like mode
\item
  How to switch environments (env var, config flag)
\item
  Hot reload behavior
\item
  Ports and URLs to access the running service
\item
  A sample invocation with expected output
\item
  How to verify the service is healthy
\end{itemize}

\noindent{}\textbf{You understand this when} you can run the project in both dev and prod modes locally and know why each mode exists.

\subsection*{3.5 Configuration Reference}\phantomsection\label{project-planning-guide:35-configuration-reference}

Every configuration variable the project supports~--- type, default, description.

\textbf{Why it matters:} Undocumented config is invisible config. Engineers who don't know an option exists won't use it correctly.

\textbf{What this topic covers:}

\begin{itemize}
\tightlist
\item
  Each variable's name, type, default
\item
  Required vs optional vs derived
\item
  What each variable controls
\item
  Format expected (URL, integer, boolean string)
\item
  Validation that runs at startup
\item
  The \texttt{.env.example} pattern for documenting variables
\end{itemize}

\noindent{}\textbf{You understand this when} you can read the configuration reference and predict the system's behavior in a new environment.

\subsection*{3.6 Testing}\phantomsection\label{project-planning-guide:36-testing}

The test strategy~--- what's tested, how, and at what level.

\textbf{Why it matters:} Untested code is owned by no one. Tests document intent, catch regressions, and enable safe change.

\textbf{What this topic covers:}

\begin{itemize}
\tightlist
\item
  The test pyramid~--- unit, integration, contract, end-to-end
\item
  How to run the full test suite
\item
  How to run a single test or test file
\item
  How to add a new test
\item
  Test doubles~--- mock, stub, fake, spy~--- when each fits
\item
  Test fixtures and data management
\item
  Coverage as a signal (its uses and limits)
\item
  Mocking strategy for external services
\item
  Flaky tests and what to do about them
\end{itemize}

\noindent{}\textbf{You understand this when} you can write a test that catches a real regression and explain why each level of the pyramid is needed.

\subsection*{3.7 Code Conventions}\phantomsection\label{project-planning-guide:37-code-conventions}

The agreed-upon style and design principles the codebase follows.

\textbf{Why it matters:} Consistency reduces cognitive load. Code that all looks the same is code that's easier to read. Code-craft principles are also the foundation under every other layer~--- design, testing, deployment, operations.

\textbf{What this topic covers:}

\begin{itemize}
\tightlist
\item
  Naming~--- files, functions, variables, constants
\item
  Linting and formatting tools
\item
  Comment style~--- when and how to comment
\item
  Module design~--- deep modules, information hiding, coupling vs cohesion
\item
  Error handling~--- how errors are raised, caught, logged; where to handle them
\item
  Function design~--- size, parameters, return shapes, side effects
\item
  Code smells~--- long methods, large classes, primitive obsession, feature envy
\item
  Refactoring as a discipline~--- small steps, under test
\item
  The cost of complexity over time
\item
  Commit message format
\end{itemize}

\noindent{}\textbf{You understand this when} you can read existing code and tell which conventions are followed, then write new code that fits in seamlessly.

\subsection*{3.8 Common Debugging Scenarios}\phantomsection\label{project-planning-guide:38-common-debugging-scenarios}

The recurring problems developers hit, and how to resolve them.

\textbf{Why it matters:} A debugging guide saves hours of repeated investigation. The ``five-minute fix you've forgotten how to do'' is exactly what this section captures.

\textbf{What this topic covers:}

\begin{itemize}
\tightlist
\item
  Per-layer diagnosis~--- where to look first (DNS, network, app, database)
\item
  Tools for each layer (strace, dig, curl, log search, journalctl, jq, database stats)
\item
  Reproducing production issues locally
\item
  Reading logs effectively~--- filtering, correlating with request IDs
\item
  Using a debugger~--- breakpoints, stepping, watching state
\item
  Common failure symptoms and their typical causes
\item
  When to escalate vs when to dig deeper
\end{itemize}

\noindent{}\textbf{You understand this when}, the moment something breaks, you can identify which layer it's in within two minutes and have a diagnostic command ready.

\unnumberedsection{true}{Section 4 --- Operational Guide}\label{project-planning-guide:section-4--operational-guide}

\subsection*{4.1 Deployment}\phantomsection\label{project-planning-guide:41-deployment}

How the application gets from source to running production~--- including build, packaging, infrastructure, and the path traffic takes to reach it.

\textbf{Why it matters:} Deployment is where code becomes a service. A team that can't deploy reliably can't ship.

\textbf{What this topic covers:}

\begin{itemize}
\tightlist
\item
  Build process~--- turning source into a deployable artifact
\item
  Packaging strategy (containers, wheels, OS packages)
\item
  Artifact storage (registries, versioning by git hash)
\item
  Infrastructure~--- servers, VMs, containers, managed services
\item
  The path from artifact to running process (pull, supervise, start)
\item
  Networking~--- DNS, reverse proxy, TLS termination, firewalls
\item
  Deployment strategies~--- recreate, rolling, blue-green, canary
\item
  Rollback procedure
\item
  How configuration and secrets are injected at runtime
\end{itemize}

\noindent{}\textbf{You understand this when} you can trace a code change from \texttt{git\ push} all the way to user traffic, naming every system it passes through.

\subsection*{4.2 Scheduling \& Triggers}\phantomsection\label{project-planning-guide:42-scheduling--triggers}

How the system handles work that isn't direct user requests~--- scheduled jobs, event-triggered tasks, webhooks.

\textbf{Why it matters:} Most production systems have hidden work running in the background. Missed runs are silent failures.

\textbf{What this topic covers:}

\begin{itemize}
\tightlist
\item
  Trigger types~--- cron, event, webhook, manual
\item
  Cron expressions and the timezone trap (cron uses the server's local time~--- schedule in UTC)
\item
  Modern alternatives to cron (systemd timers, cloud schedulers)
\item
  Event-triggered work~--- webhooks, message queues
\item
  Overlap handling~--- what if one run is still active when the next fires
\item
  Idempotency for scheduled jobs~--- essential; jobs must be safely re-runnable
\item
  Durability~--- what if the scheduler or job crashes mid-execution
\item
  Monitoring missed runs (not just failed runs)
\end{itemize}

\noindent{}\textbf{You understand this when} for every scheduled job in the system, you can describe what triggers it, what happens if it fails, and how you'd notice if it stopped running.

\subsection*{4.3 Monitoring \& Health Checks}\phantomsection\label{project-planning-guide:43-monitoring--health-checks}

How you know the system is alive and behaving correctly.

\textbf{Why it matters:} A system you can't observe is a system you can't operate. Problems must surface before users notice.

\textbf{What this topic covers:}

\begin{itemize}
\tightlist
\item
  The three pillars~--- logs, metrics, traces
\item
  Structured logging (JSON) and request ID correlation
\item
  Health check endpoints~--- deep vs shallow
\item
  The four golden signals~--- latency, traffic, errors, saturation
\item
  Metrics~--- counters, gauges, histograms
\item
  Log aggregation
\item
  Alerting on symptoms (not causes), and the alert-fatigue trap
\item
  SLIs, SLOs, error budgets
\end{itemize}

\noindent{}\textbf{You understand this when} you can answer ``is the system healthy right now?'' with a single command, and you trust the answer.

\subsection*{4.4 Common Failure Modes}\phantomsection\label{project-planning-guide:44-common-failure-modes}

The known ways the system can fail in production, and how to diagnose each one.

\textbf{Why it matters:} Engineers diagnose faster when failures are familiar. A documented failure mode is a failure that's already half-solved.

\textbf{What this topic covers:}

\begin{itemize}
\tightlist
\item
  Per-layer failures~--- server (OOM, disk full), reverse proxy (cert expiry, misconfig), container (crash loop, unhealthy), application (worker timeout, queue saturation), database (connection refused, slow queries)
\item
  Symptoms~--- what the user sees
\item
  Detection~--- what to look for in logs/metrics
\item
  Diagnosis commands~--- what to run to confirm
\item
  Recovery steps
\item
  Prevention measures for the future
\end{itemize}

\noindent{}\textbf{You understand this when} you can pattern-match an alert to a likely failure mode in under a minute and start diagnosis with the right command.

\subsection*{4.5 Restart \& Recovery Procedures}\phantomsection\label{project-planning-guide:45-restart--recovery-procedures}

The step-by-step procedure for restoring a service after a failure.

\textbf{Why it matters:} Under stress, fresh thinking fails. A documented runbook is the difference between five-minute and five-hour recovery.

\textbf{What this topic covers:}

\begin{itemize}
\tightlist
\item
  Assess~--- how to determine the scope of the failure
\item
  Stop~--- how to cleanly stop a running process if needed
\item
  Fix~--- what commonly needs to be resolved before restart
\item
  Restart~--- the exact command
\item
  Verify~--- how to confirm the system is healthy
\item
  Process supervisors and restart policies
\item
  Graceful shutdown~--- draining in-flight requests, signal handling
\item
  The PID 1 signal-handling gotcha in containers
\item
  Restoring from backup during an incident (strategy and restore testing live in 4.6)
\item
  Postmortem documentation and blameless culture
\end{itemize}

\noindent{}\textbf{You understand this when} you can recover the service from a cold start without consulting documentation~--- but you wrote the runbook anyway.

\subsection*{4.6 Data Management}\phantomsection\label{project-planning-guide:46-data-management}

How data is retained, archived, deleted, and protected over time.

\textbf{Why it matters:} Real systems accumulate data that's expensive to keep and legally fraught to delete badly. Both retention and deletion are increasingly regulated.

\textbf{What this topic covers:}

\begin{itemize}
\tightlist
\item
  Retention policy~--- how long different data lives
\item
  Storage tiers~--- hot, warm, cold, archival
\item
  Archival procedures~--- moving old data without losing it
\item
  Cleanup jobs~--- safe deletion with dry-runs and audit logs
\item
  PII handling~--- minimization, pseudonymization, what to log and what not to
\item
  Backup strategy~--- frequency, storage, verification
\item
  Restore testing~--- an untested backup is a hope, not a backup
\item
  Compliance-driven retention (GDPR right-to-delete, legal hold)
\item
  The difference between soft delete, hard delete, and cryptographic deletion
\item
  Local vs production parity in data behavior
\end{itemize}

\noindent{}\textbf{You understand this when} for every category of data in the system, you can describe how long it lives, how it's backed up, how it would be deleted, and how a data-subject deletion request would be handled.

\unnumberedsection{true}{Section 5 --- Delivery \& Changelog}\label{project-planning-guide:section-5--delivery--changelog}

\subsection*{5.1 Delivery Phases}\phantomsection\label{project-planning-guide:51-delivery-phases}

The discrete milestones the project is broken into~--- each with a clear goal and definition of done.

\textbf{Why it matters:} Without phases, projects either drift indefinitely or get rushed at the end. Phases create natural checkpoints.

\textbf{What this topic covers:}

\begin{itemize}
\tightlist
\item
  Each phase's goal
\item
  Scope included in each phase
\item
  Status (completed / in progress / planned)
\item
  Target or actual completion date
\item
  The distinction between phases (stable milestones) and tasks (individual work items)
\item
  How to handle scope additions during a phase
\end{itemize}

\noindent{}\textbf{You understand this when} you can describe the project's phases to a stakeholder and predict when each one delivers value.

\subsection*{5.2 Current State}\phantomsection\label{project-planning-guide:52-current-state}

A snapshot of what's live in production right now.

\textbf{Why it matters:} Without a current-state record, no one knows what version is deployed or what works. New team members can't get up to speed.

\textbf{What this topic covers:}

\begin{itemize}
\tightlist
\item
  What features are live in production
\item
  Current version (release tag or git hash)
\item
  Date of most recent change to production
\item
  Known limitations of the current state
\item
  Active workarounds in place
\item
  The relationship between current state and the roadmap
\end{itemize}

\noindent{}\textbf{You understand this when} any team member can answer ``what's in production right now?''~--- and the answer matches what the deployment system reports.

\subsection*{5.3 Changelog}\phantomsection\label{project-planning-guide:53-changelog}

A running record of significant changes~--- what was changed, when, and by whom.

\textbf{Why it matters:} Without a changelog, releases are opaque. Users can't tell what's new; engineers can't tell what's broken.

\textbf{What this topic covers:}

\begin{itemize}
\tightlist
\item
  Date of each significant change
\item
  Version number (semantic versioning~--- major.minor.patch)
\item
  Author
\item
  Category of change~--- added, changed, deprecated, removed, fixed, security
\item
  Conventions~--- Keep a Changelog, Conventional Commits
\item
  The difference between deploys (technical) and releases (user-visible)
\item
  Release notes for stakeholders vs developer changelog
\end{itemize}

\noindent{}\textbf{You understand this when} you can read the changelog and reconstruct the project's history at a glance.

\subsection*{5.4 Known Issues \& Tech Debt}\phantomsection\label{project-planning-guide:54-known-issues--tech-debt}

The acknowledged problems that haven't been fixed yet~--- and the trade-offs made deliberately.

\textbf{Why it matters:} Tech debt nobody tracks compounds silently. Acknowledged debt can be paid down strategically.

\textbf{What this topic covers:}

\begin{itemize}
\tightlist
\item
  Each known issue with impact~--- who or what is affected
\item
  The current workaround
\item
  The proper fix
\item
  Each piece of tech debt with: why it was done this way, what it would take to clean up
\item
  The technical-debt quadrant (Martin Fowler)~--- deliberate vs inadvertent, prudent vs reckless
\item
  How to prioritize debt against new feature work
\end{itemize}

\noindent{}\textbf{You understand this when} you can defend the prioritization between paying down a piece of debt and shipping a new feature.

\subsection*{5.5 Future Roadmap}\phantomsection\label{project-planning-guide:55-future-roadmap}

What comes next~--- features and improvements planned, prioritized.

\textbf{Why it matters:} A roadmap aligns the team and stakeholders on direction. Without one, every priority debate happens in isolation.

\textbf{What this topic covers:}

\begin{itemize}
\tightlist
\item
  Each upcoming feature or improvement
\item
  Priority (High / Medium / Low)
\item
  Rationale for each priority
\item
  Dependencies between roadmap items
\item
  Time horizon (next quarter, next year, eventually)
\item
  How the roadmap is revised when reality intervenes
\end{itemize}

\noindent{}\textbf{You understand this when} you can explain to a new team member what's coming in the next quarter and why.

\unnumberedsection{true}{A Worked Example: Planning the Notes API}\label{project-planning-guide:a-worked-example-planning-the-notes-api}

The template is easier to use once you have seen it filled in. Here is a compressed plan for the notes service whose discovery was sketched in Part I~--- the same Notes API that Part IV builds, chapter by chapter. A real plan would say more in each section; the point here is what \emph{one good line} per section looks like.

\textbf{Section 1~--- Problem Statement \& Goals}

\begin{itemize}
\tightlist
\item
  \textbf{Background \& context:} discovery showed that developers lose the snippets and decisions they jot down during the day; a two-week trial with a shared file confirmed they would use a simple capture tool daily.
\item
  \textbf{Core problem:} capturing a note must take one command from wherever the developer already is, and finding it again must be just as quick.
\item
  \textbf{Goals \& success criteria:} a note can be created, listed, read, updated, and deleted over HTTPS; success is ``each pilot user retrieves a saved note at least once a week'' after one month.
\item
  \textbf{Out of scope (v1):} user accounts and sharing, a web or mobile front end, full-text search ranking, attachments.
\item
  \textbf{Stakeholders \& users:} the twelve pilot developers (users); the one engineer building it (owner); the team lead who pays the hosting bill.
\item
  \textbf{Constraints \& assumptions:} one engineer, six weeks, hosting under \$25 a month; assumes a single region and a few requests per second at most.
\end{itemize}

\noindent{}\textbf{Section 2~--- Architecture \& Design}

\begin{itemize}
\tightlist
\item
  \textbf{System overview:} a small REST API~--- Python and Flask~--- in a Docker container on one Linux VM, behind Nginx for TLS, with a managed PostgreSQL database.
\item
  \textbf{Architecture diagram:} client → DNS → cloud firewall → Nginx (TLS) → Gunicorn in Docker → managed PostgreSQL over a private network.
\item
  \textbf{Component breakdown:} application factory, configuration module, database access, route handlers, validation, error handlers, logging.
\item
  \textbf{Data flow:} a \texttt{POST\ /notes/} is validated, inserted, committed, and returned as JSON; every request is logged with its status and duration.
\item
  \textbf{External integrations:} none in v1 beyond the database~--- which is still a contract, with its own connection string, credentials, and failure modes.
\item
  \textbf{Data storage \& schema:} one table, \texttt{notes(id,}\allowbreak{}\texttt{\ content,}\allowbreak{}\texttt{\ created\_}\allowbreak{}\texttt{at,}\allowbreak{}\texttt{\ updated\_}\allowbreak{}\texttt{at)}; schema changes go through versioned migrations.
\item
  \textbf{Key design decisions:} SQLite for local development and PostgreSQL in production, recorded as an ADR~--- and later superseded by ``PostgreSQL everywhere'' when the difference hid a bug.
\item
  \textbf{Security considerations:} HTTPS only; all input validated; secrets in environment files, never in the repository; no authentication in v1 is an \emph{accepted, written-down} risk: the pilot's URL is shared only with the pilot users, the data is low-stakes, and authentication is first on the roadmap.
\end{itemize}

\noindent{}\textbf{Section 3~--- Developer Guide}

\begin{itemize}
\tightlist
\item
  \textbf{Prerequisites:} Python 3.11, Docker, git.
\item
  \textbf{Repository structure:} \texttt{app/} (the package), \texttt{tests/}, \texttt{migrations/}, \texttt{Dockerfile}, \texttt{Makefile}, \texttt{requirements*.txt}.
\item
  \textbf{Environment setup:} create a virtual environment, \texttt{pip\ install\ -}\allowbreak{}\texttt{r\ requirements-}\allowbreak{}\texttt{dev.}\allowbreak{}\texttt{txt}, copy \texttt{.env.example} to \texttt{.env} and fill in \texttt{SECRET\_KEY}.
\item
  \textbf{Running locally:} \texttt{python\ run.py}, then \texttt{curl\ http:}\allowbreak{}\texttt{/}\allowbreak{}\texttt{/}\allowbreak{}\texttt{localhost:}\allowbreak{}\texttt{5000/}\allowbreak{}\texttt{health}.
\item
  \textbf{Configuration reference:} \texttt{DATABASE\_URL} and \texttt{SECRET\_KEY} (required); \texttt{FLASK\_DEBUG}, \texttt{PORT}, \texttt{HOST}, \texttt{LOG\_LEVEL}, \texttt{LOG\_FORMAT} (optional, with defaults).
\item
  \textbf{Testing:} \texttt{make\ build} runs lint, type checks, the test suite, and a dependency audit; CI runs the same on every push.
\item
  \textbf{Code conventions:} formatted by black, linted by ruff; one error-response shape for the whole API.
\item
  \textbf{Common debugging scenarios:} ``port already in use'', a missing environment variable at startup, a 415 because a client forgot \texttt{Content-}\allowbreak{}\texttt{Type:}\allowbreak{}\texttt{\ application/}\allowbreak{}\texttt{json}.
\end{itemize}

\noindent{}\textbf{Section 4~--- Operational Guide}

\begin{itemize}
\tightlist
\item
  \textbf{Deployment:} CI builds and smoke-tests an image tagged by commit, pushes it to the registry, and runs \texttt{deploy.sh\ sha-\textless{}commit\textgreater{}} on the server, which rolls back automatically if the health check fails.
\item
  \textbf{Scheduling \& triggers:} none in v1; a nightly cleanup job would be the first.
\item
  \textbf{Monitoring \& health checks:} \texttt{/health} checks the database and is watched by an external uptime monitor; structured JSON logs; alerts on downtime and error rate.
\item
  \textbf{Common failure modes:} database unreachable (503 from \texttt{/health}), expired TLS certificate, disk full from logs, a crash-looping container after a bad deploy.
\item
  \textbf{Restart \& recovery:} redeploy the last good tag; the service-down runbook covers the first ten minutes of any outage.
\item
  \textbf{Data management:} managed-database backups kept 7 days, plus an off-server copy before every schema change; a monthly restore test.
\end{itemize}

\noindent{}\textbf{Section 5~--- Delivery \& Changelog}

\begin{itemize}
\tightlist
\item
  \textbf{Delivery phases:} (1) runs locally, (2) portable and configurable, (3) packaged as an image, (4) deployed to a server, (5) reachable over HTTPS, (6) observable and hardened~--- each a working system before the next begins.
\item
  \textbf{Current state:} v1.2.0 in production (image \texttt{sha-a1b2c3d}), deployed on the date in the changelog.
\item
  \textbf{Changelog:} kept in \texttt{CHANGELOG.md} under \emph{Added / Changed / Fixed / Security}, one entry per release.
\item
  \textbf{Known issues \& tech debt:} no authentication (accepted for the pilot); one server is a single point of failure; the list endpoint needs pagination before the data grows.
\item
  \textbf{Future roadmap:} High~--- authentication, before anyone outside the pilot gets the URL; Medium~--- pagination and search; Low~--- a small web front end.
\end{itemize}

\noindent{}Filled in like this, the plan fits on three pages, and every line is something a newcomer, an operator, or a future you will need to know.

\unnumberedsection{true}{How This Guide Evolves}\label{project-planning-guide:how-this-guide-evolves}

This Guide is a living project artifact. Complete each section before or during early development, and return to update it as the project evolves~--- when architecture decisions are made, when scope changes, when new failure modes are discovered.

Over time, the \textbf{Why it matters} and \textbf{What this topic covers} prompts recede into the background and the sections fill with real project content. The Guide is alive. Each decision made, each component added, and each incident resolved is an input to it.
